\chapter{完成情况及性能参数}

\section{整体介绍}

{[}待填写：工程版本与源码提交、目标平台与器件、运行配置、工具版本、整体完成情况，以及仿真验证、FPGA实现和实物验证的范围与边界。{]}

\section{工程成果}

\subsection{硬件成果}

{[}待填写：CPU、NPU、系统互连、存储及图像采集与显示链路的实现成果，FPGA集成情况，以及对应RTL、板级工程、约束文件和硬件交付清单。{]}

\subsection{软件成果}

{[}待填写：模型导出与量化部署、裸机运行时、外设和NPU驱动、CPU回退路径、运行程序及配套工具的实现成果，以及软件交付清单和复现入口。{]}

\subsection{功能验证}

{[}待填写：CPU指令执行、NPU计算、数据搬运、共享存储访问、图像链路和完整模型运行的验证结论，已通过项目、待验证项目与已知限制，以及对应附录中的测试用例、波形和日志索引。{]}

\subsection{应用演示}

{[}待填写：系统连接与板卡接口、启动和操作流程、输入输出图像、连续帧运行效果及演示条件；分别标明仿真演示与实物演示，并列出实物照片、传感器、显示输出和运行画面的实际验证状态。{]}

\section{特性成果}

\subsection{CPU性能}

{[}待填写：测试程序、编译选项与运行配置、执行周期、退休指令数、IPC及相关性能计数器，统计口径、测试结果和影响因素；区分仿真结果与板上实测结果，注明基准测试的运行条件和结果适用范围。{]}

\subsection{NPU性能}

{[}待填写：测试算子与矩阵规模、计算周期、DMA数据搬运开销、阵列利用率，以及相同任务和数值约定下相对CPU执行的加速比；说明计时范围、运行配置和测量方法，区分计算单元性能与包含搬运的任务性能。{]}

\subsection{系统性能}

{[}待填写：模型与输入尺寸、完整模型推理延迟、端到端处理延迟、连续帧吞吐率、CPU与NPU协同执行开销，以及系统瓶颈分析；说明各项指标的计时起止点、测试条件和证据类型，区分仿真推算与板上实测。{]}

\subsection{FPGA实现指标}

{[}待填写：目标器件与工程配置、逻辑资源、寄存器、DSP和存储资源占用，工作频率、时序裕量及关键路径；注明报告来自综合还是布局布线阶段，并区分时序约束目标、实现结果和实际板上工作条件。{]}
